\documentclass[11pt,a4paper]{article}
\usepackage[T1]{fontenc}
\usepackage{lmodern,amsmath,amssymb,booktabs,microtype}
\usepackage[margin=24mm]{geometry}
\usepackage[colorlinks=true,urlcolor=blue,linkcolor=blue]{hyperref}
\newcommand{\R}{\mathbb R}
\newcommand{\T}{\mathbb T}
\newcommand{\Z}{\mathbb Z}
\newcommand{\F}{F_\hbar}
\newcommand{\dd}{\,d}
\title{Paired quantum dilogarithms:\\source dictionary and bounded kernel tests}
\author{Alan Day \and Chengheng\thanks{AI research assistant (OpenAI).}}
\date{Working research note --- 1 November 2026}
\begin{document}\maketitle
\begin{abstract}
Existing Gaussian-group Fourier identities offer a direct language for
the bounded paired pentagon. We match their measure and sign conventions,
distinguish three nearby constructions, and give a heat-regularized
kernel test without unbounded chiral factors. Its matrix-element error
reduces to a positive Gaussian phase-loss expectation. An empirical
Sobol study explores the size of that functional, but does not certify
it or evaluate the pentagon defect. The original bounded identity and
the physical observer gluing problem remain open in this project.
\end{abstract}
\section{Literature first: what can be imported}
The scalar function under study is the Kim--Scarinci (KS) phase
\cite{ks}
\[
 \F(x,y)=
 \frac{\Phi^{i\hbar}(x+i\hbar y)}
 {\overline{\Phi^{i\hbar}(x+i\hbar y)}},\qquad \hbar>0 .
\]
It has modulus one almost everywhere, with arbitrary values assigned
at the discrete exceptional points. Coordinates are dimensionless.
The claim being reproduced is KS Proposition~6.15, in the representation
\[
 Q=(X,Y),\quad P=(P_y,P_x),\quad P_j=-\pi i\partial_j,\quad
 A=\F(Q),\quad B=\F(P),\quad M=e^{iXY/\pi},
\]
\[
 AB=BM^{-1}BMA. \tag{1}
\]
The components of each vector strongly commute. The shift in the middle
factor is the self-adjoint closure specified by
$M^{-1}PM=P+Q$.

Garoufalidis--Kashaev (GK), Appendix~B, equations (136)--(138)
\cite{index}, provides the bounded-operator/Fourier-kernel equivalence
on a Gaussian group. GK's later Remark~2.6 restates the distributional
version \cite{local}. These results give a framework, not an automatic
proof that our particular function satisfies its required identity.
Kashaev's earlier Lemma~1 and Theorem~1 have square-integrability
hypotheses \cite{beta}; a unit-modulus function on $\R^2$ is not in
$L^2(\R^2)$. A damping factor supplies integrability but generally
destroys the exact pentagon, so it cannot be inserted into those
theorems as though the hypothesis were unchanged.

\newpage
\section{The precise Gaussian group and nearby examples}
Use the additive group $G=\R^2$ with
\[
 \gamma(s)=e^{-is_xs_y/\pi},\qquad
 \chi(s,t)=e^{-i(s_xt_y+s_yt_x)/\pi},\qquad
 d\mu(s)=\frac{ds_x\,ds_y}{2\pi^2}.
\]
The negative-sign Fourier transform $S$ with kernel $\chi$ is unitary,
$S^2f(s)=f(-s)$, and its Gaussian integral is $1$ in the oscillatory
sense. Thus $SQS^*=-P$. The source convention is
$f(q)=f(Q)$, $f(p)=Sf(Q)S^*=f(-P)$ and
$f(p+q)=\gamma(Q)^*f(p)\gamma(Q)=f(Q-P)$.
Conjugating the source pentagon by $S^*$ gives exactly (1).

With $k=S^*\F$ as a tempered distribution, the matched version
of GK's equation (137) is
\[
 k(a)k(b)\chi(a,b)=
 \int_G k(b-c)k(c)k(a-c)\gamma(c)\,d\mu(c). \tag{2}
\]
Equation (2) is the identity to establish, not an ordinary absolutely
convergent integral already known for this $k$.

\paragraph{The known real-plane example.}
GK Example~2.4 uses
$\varphi_b(\dot x+\ddot x/2)/\varphi_b(\dot x-\ddot x/2)$, with the
usual Faddeev convention. For real $b=\sqrt\hbar$, set
\[
 \dot x=\frac{x}{2\pi\sqrt\hbar},\qquad
 \ddot x=-\frac{\sqrt\hbar\,y}{\pi}.
\]
The absolute Jacobian is $1/(2\pi^2)$ and
$e^{2\pi i\dot x\ddot x}=\gamma(x,y)$.
Since the KS convention is
$\Phi^\hbar(z)=\varphi_{\sqrt\hbar}(z/(2\pi\sqrt\hbar))^{-1}$,
this gives the KS \emph{negative}-cosmological-constant pairing
\[
 \frac{\Phi^\hbar(x+\hbar y)}{\Phi^\hbar(x-\hbar y)}.
\]
Replacing $\hbar$ by $i\hbar$ is not a real change of Haar coordinates.
The real-parameter ratio is smooth on the real plane, whereas our
phase has isolated nonremovable windings. Thus a real invertible
linear substitution, even with a smooth nonvanishing phase, cannot
identify these two multipliers. This does not exclude a more general
unitary equivalence.

\paragraph{The local-field construction.}
GK Definition~3.1 and Theorem~3.3 use $B=F^\times$,
$G=\widehat B\times B$, and $\varphi(\alpha,z)=\alpha(1+z)$.
For $F=\mathbb C$ this group is $\R^2\times\Z\times\T$, not
the additive complex plane. Lemma~3.2 gives a Fourier expression
away from $z=1$ and explicitly retains its distributional qualification
there. The theorem does not identify this construction with $\F$.

\newpage
\section{A scalar cylinder dictionary is not an operator dictionary}
For real $0<q<1$, write
\[
 C_q(\theta,m)=
 \frac{(-q^{1-m}e^{i\theta};q^2)_\infty}
      {(-q^{1-m}e^{-i\theta};q^2)_\infty}.
\]
For integer $m$, this is the $\T\times\Z$ phase in GK
Appendix~A, equations (129)--(131) \cite{index}. Away from exceptional
arguments, the same product is a scalar function of real $m$.
The KS product formula gives the exact scalar factorization
\[
 \F(x,y)=
 \frac{C_{\widetilde q}(-x/\hbar,\hbar y/\pi)}
      {C_q(\hbar y,x/(\pi\hbar))},\qquad
 q=e^{-\pi\hbar},\quad\widetilde q=e^{-\pi/\hbar}. \tag{3}
\]
Both second arguments in (3) are generally noninteger.

The known inversion formula
\[
 C_q(\theta,m)C_q(-\theta,-m)=e^{im\theta},\qquad m\in\Z,
\]
cannot hold for arbitrary real $m$ as a function on the circle.
The left side is $2\pi$-periodic, whereas the proposed right side
has monodromy $e^{2\pi im}$. Passing to the real universal cover
changes the group and its Fourier dual; it does not extend the
original operator theorem without further work. Likewise, integer
radial spectra in the quantum-$az+b$ theory cannot silently be
replaced by the full complex spectrum of $e^{X+i\hbar Y}$.

The \emph{combined} phase in (3) still has the correct inversion,
\[
 \F(x,y)\F(-x,-y)=e^{-ixy/\pi}. \tag{4}
\]
It follows directly from KS's scalar inversion formula, including
its real constant at parameter $i\hbar$. Thus the individual
continuous-index failure does not prove a failure of the paired
pentagon. Sixty-digit checks of nine generic $(\hbar,x,y)$ points
verify (3) and (4); four integer cylinder controls and three
noninteger monodromy controls are also saved.

\paragraph{What the search has established.}
There is an applicable \emph{formulation} of the bounded identity.
The real-plane, cylinder, and local-field examples do not yet
supply its missing proof for our phase in this representation.
These are identified hypothesis differences, not a claim that
all relevant literature has been exhausted. We do not claim a
new general quantum-dilogarithm theorem.

\newpage
\section{An absolutely defined regulated kernel}
Let
\[
 H_\epsilon=e^{-\epsilon(P_x^2+P_y^2)},\quad
 B_\epsilon=BH_\epsilon,\quad
 k_\epsilon=S^*(e^{-\epsilon|s|^2}\F(s)),\qquad\epsilon>0.
\]
The Fourier integral defining $k_\epsilon$ is absolutely convergent,
and Plancherel gives
\[
 \|k_\epsilon\|_\infty\leq\frac1{2\pi\epsilon},
 \qquad
 \|k_\epsilon\|_{L^2(\mu)}^2=\frac1{4\pi\epsilon}.
\]
The convolution kernel of
$D_\epsilon=AB_\epsilon-B_\epsilon M^{-1}B_\epsilon MA$ is therefore
\begin{align*}
 D_\epsilon(s,t)
 &=\F(s)k_\epsilon(s-t)\\
 &\quad-M(t)\F(t)\int_G
 k_\epsilon(s-u)\overline{M(u)}k_\epsilon(u-t)\,d\mu(u).
\end{align*}
The inner integral is absolutely convergent by Cauchy--Schwarz,
bounded in modulus by $1/(4\pi\epsilon)$. For Schwartz input and
output tests the remaining double pairing is also absolutely
convergent. This does not authorize expanding every Fourier
integral and exchanging all oscillatory integrations at once.

Dominated convergence in the momentum spectral measure gives
$B_\epsilon\to B$ strongly. Products are uniformly bounded, so
$D_\epsilon\to D=AB-BM^{-1}BMA$ strongly. No unbounded chiral
factor or fixed-vector holomorphic continuation is used.

For normalized $f,g$, put $v=MAf$ and
$E_\epsilon(w)=\|(1-H_\epsilon)w\|$. A telescoping difference and
Cauchy--Schwarz yield the useful matrix-element bound
\[
 |\langle g,(D-D_\epsilon)f\rangle|
 \leq E_\epsilon(f)+E_\epsilon(g)+E_\epsilon(v). \tag{5}
\]
For example, moving the outer $B-B_\epsilon$ to the output test
costs $E_\epsilon(g)$, since $B$ commutes with $H_\epsilon$ and is
unitary. The inner difference costs $E_\epsilon(MAf)$.
All three errors tend to zero for every fixed $L^2$ test.
One must still compute and control a matrix element of $D_\epsilon$;
its value need not be zero even if $D=0$.

\newpage
\section{Gaussian reduction of the only rough heat tail}
In $L^2(\mu)$ take $f=g=f_0=\sqrt2\,e^{-|s|^2/(2\pi)}$,
so $\|f_0\|=1$. Then
\[
 E_\epsilon(f_0)^2
 =\frac{2(\pi\epsilon)^2}{(1+\pi\epsilon)(1+2\pi\epsilon)}.
\]
For $v(s)=r(s)f_0(s)$, $r=M\F$, define the real heat overlap
$C_t=\langle v,H_t v\rangle$. In ordinary coordinates the heat
kernel is $(4\pi^3t)^{-1}e^{-|s-s'|^2/(4\pi^2t)}$.
Using $U=(s+s')/2$, $Z=s-s'$, Gaussian completion gives
\[
 C_t=\frac1{1+\pi t}\,
 \mathbb E\!\left[
 \overline{r(U+Z/2)}r(U-Z/2)\right], \tag{6}
\]
where $U,Z$ are independent centered two-dimensional Gaussians with
\[
 \operatorname{Cov}(U)=\frac\pi2 I,\qquad
 \operatorname{Cov}(Z)=\frac{2\pi^2t}{1+\pi t}I .
\]
Symmetry makes the expectation real. An explicitly nonnegative
form, useful near cancellation, is
\[
 1-C_t=\frac{\pi t+\mathbb E L_t}{1+\pi t},\qquad
 L_t=\frac12|r(U+Z/2)-r(U-Z/2)|^2\in[0,2]. \tag{7}
\]
The spectral inequality $(1-a)^2\leq1-a^2$, $0\leq a\leq1$, proves
\[
 E_\epsilon(v)^2\leq1-C_{2\epsilon}.
\]
Combining it with (5) yields the exact bound functional
\[
 {\cal E}_\epsilon=
 2\sqrt{\frac{2(\pi\epsilon)^2}
 {(1+\pi\epsilon)(1+2\pi\epsilon)}}+\sqrt{1-C_{2\epsilon}}. \tag{8}
\]
No $H^1$ assumption on the vortex phase is needed.
A numerical estimate of (7) is not itself a certified upper bound.

As a separate analytic control, remove $F$ but retain $r=M$.
Integrating the Gaussian bilinear phase in (6) gives
$C_t=(1+2\pi t)^{-1}$. If $r=1$, instead
$C_t=(1+\pi t)^{-1}$, which recovers the exact vacuum formula above.
These controls test the measure, covariance and chirp normalization.

\newpage
\section{Exploration, evidence and the next decision}
We used four scrambled Sobol seeds, $2^{16}$ four-dimensional
normal samples per seed, three paired cases
$\hbar=.5,1,2$, the analytic chirp control, and five
$\epsilon$ values from $10^{-2}$ to $10^{-6}$.
Prefixes $2^{12},2^{14},2^{16}$ give 240 saved records.
Prefixes and settings reuse base nodes; they are not independent
replications. Seeds, versions, input hashes, prefix moments and
counts of large phase losses are saved.

For $\hbar=1$, the four fine-grid estimates of (8) span:
\begin{center}
\begin{tabular}{rcc}
\toprule
$\epsilon$ & minimum & maximum\\
\midrule
$10^{-2}$ & .376212 & .376429\\
$10^{-3}$ & .106305 & .106500\\
$10^{-4}$ & .0319381 & .0324646\\
$10^{-5}$ & .00991042 & .0101192\\
$10^{-6}$ & .00311482 & .00318232\\
\bottomrule
\end{tabular}
\end{center}
At $10^{-6}$ the ranges for $\hbar=.5,2$ are approximately
$.00307294$--$.00307389$ and $.00307308$--$.00307398$.
The chirp-only control has maximum relative deviation
$2.80\,10^{-4}$ at the fine order. Thirty scalar, symbolic and
control checks pass. These figures explore feasibility only.
The larger replicate spread at $\hbar=1$ is a reason to resolve
rare phase-loss regions, not evidence of a certified error bar.

\paragraph{Next scientific calculation.}
Use the matched bounded kernel to compute one specified
$\langle f_0,D_\epsilon f_0\rangle$ with certified integration error,
and separately bound the expectation in (7). Split off Gaussian
tails and neighborhoods of the known phase singularities; their
contribution is controlled by $0\leq L_t\leq2$.
Only a certified modulus exceeding the sum of those errors can
yield a continuous nonzero defect. If it does not, retain that
outcome and refine the representation/identity comparison.
The earlier finite-grid residuals and the present empirical
heat estimates cannot be combined into a counterexample.

This supplies an executable route around the unbounded-factor
difficulty, using known Fourier and heat-semigroup facts.
It does not resolve the final pentagon, provide observer mass labels
or construct the finite-$G$ physical gluing measure.

\begin{thebibliography}{9}\small
\bibitem{ks} H.~K.~Kim and C.~Scarinci,
\emph{A quantization of moduli spaces of 3-dimensional gravity},
\href{https://arxiv.org/abs/2112.13329v2}{2112.13329v2},
\S6.3, Proposition~6.15.
\bibitem{index} S.~Garoufalidis and R.~Kashaev,
\emph{A meromorphic extension of the 3D index},
\href{https://arxiv.org/abs/1706.08132v2}{1706.08132v2},
Appendices A--B, equations (129)--(138).
\bibitem{local} S.~Garoufalidis and R.~Kashaev,
\emph{Quantum dilogarithms over local fields and invariants of 3-manifolds},
\href{https://arxiv.org/abs/2306.01331v2}{2306.01331v2},
Example~2.4, Remark~2.6, Definition~3.1, Lemma~3.2, Theorem~3.3.
\bibitem{beta} R.~Kashaev, \emph{On beta pentagon relations},
\href{https://arxiv.org/abs/1403.1298v1}{1403.1298v1},
Lemma~1 and Theorem~1.
\end{thebibliography}
\end{document}
